\documentclass[aspectratio=169,11pt]{beamer}

\usetheme{default}
\usecolortheme{seahorse}
\setbeamertemplate{navigation symbols}{}
\setbeamertemplate{footline}[frame number]

\usepackage[T1]{fontenc}
\usepackage{lmodern}
\usepackage{amsmath,amssymb}
\usepackage{booktabs}

\newcommand{\HT}{\mathrm{HT}}
\newcommand{\HTp}{\mathrm{HT}^{+}}
\newcommand{\Otil}{\widetilde{O}}

\title{Quantum Walk Search as a Model of Cognitive Search}
\subtitle{From Marked Vertices to Quantum Cognition}
\author{[Student Name] \quad [Student ID]}
\institute{EE5105 Special Topics in Quantum Information Science\\Prof. Shih-Han Hung}
\date{Fall 2026}

\begin{document}

\begin{frame}
  \titlepage
\end{frame}

\begin{frame}{Takeaway}
  \begin{itemize}
    \item QIS core: AGJK finds marked vertices on any graph in $\Otil(\sqrt{\HT})$ time.
    \item Cognitive mapping: marked vertices are target memories or decision alternatives.
    \item Question: how do topology, marked-set structure, and decoherence change the advantage?
  \end{itemize}
\end{frame}

\begin{frame}{Outline}
  \tableofcontents
\end{frame}

\section{Motivation and Questions}

\begin{frame}{Why quantum walks for cognition?}
  \begin{itemize}
    \item Quantum cognition uses quantum probability without claiming a quantum brain.
    \item Random walks model evidence accumulation and semantic spreading.
    \item Quantum walks add interference between search paths.
    \item Gap: algorithmic structure of marked-vertex search is rarely mapped to cognition.
  \end{itemize}
\end{frame}

\begin{frame}{Research questions}
  \begin{enumerate}
    \item How do $\HT$, $\HTp$, spectral gap, and success probability map to response time, choice probability, and deliberation speed?
    \item How do single, independent, and clustered marked sets change the quantum advantage?
    \item Does partial measurement create an optimal level of introspection?
  \end{enumerate}
\end{frame}

\section{Background}

\begin{frame}{Classical walk and hitting time}
  \begin{itemize}
    \item Graph $G=(V,E)$, transition matrix $P$, marked set $M\subset V$.
    \item Hitting time $\HT=\mathbb{E}_{v\sim\pi}[\min\{t:X_t\in M\}]$.
    \item Examples: complete graph $\Theta(N)$, cycle $\Theta(N^2)$, grid $\Theta(N\log N)$.
  \end{itemize}
\end{frame}

\begin{frame}{Szegedy walk: detection is not finding}
  \begin{itemize}
    \item $V(P)|0\rangle|x\rangle=\sum_y\sqrt{P_{xy}}|y,x\rangle$.
    \item $W(P)=V^\dagger(P)\,\mathrm{Shift}\,V(P)\,\mathrm{Ref}$.
    \item Discriminant matrix $D_{xy}=\sqrt{P_{xy}P_{yx}}$.
    \item Detection takes $O(\sqrt{\HT})$, but measuring the final state need not output a marked vertex.
  \end{itemize}
\end{frame}

\section{AGJK}

\begin{frame}{Extended hitting time and interpolated walks}
  \begin{itemize}
    \item Absorbing walk $P'=\begin{pmatrix}P_{UU}&P_{UM}\\0&I\end{pmatrix}$.
    \item Interpolated walk $P(s)=(1-s)P+sP'$.
    \item $\HTp=\lim_{s\to1}\HT(s)$ can be much larger than $\HT$ for multiple marked vertices.
    \item AGJK: run the interpolated walk, or combine it with quantum fast-forwarding.
  \end{itemize}
\end{frame}

\section{Cognitive Model}

\begin{frame}{Graph to cognition}
  \begin{itemize}
    \item Vertices: concepts, evidence states, or options.
    \item Edges: semantic associations or evidence transitions.
    \item Marked set: target memory or decision boundary.
    \item Success probability $p_M(t)$: target activation; first threshold crossing: response time.
    \item Multiple marked sets: multi-alternative choice probabilities.
  \end{itemize}
\end{frame}

\section{Methods}

\begin{frame}{Experimental setup}
  \begin{itemize}
    \item Graphs: complete, cycle, 2D grid, hypercube, random regular, barbell.
    \item Semantic graph: hand-built 19-node toy network.
    \item Decision graphs: line and multi-branch.
    \item Marked sets: single, independent, clustered, option sets.
    \item Metrics: $T_q$, $p_{\max}$, $T_c$, speedup, choice probability, response time.
  \end{itemize}
\end{frame}

\section{Results}

\begin{frame}{Scaling}
  \begin{itemize}
    \item Classical versus quantum search time on at least two graph families.
    \item Fit $T\sim N^\alpha$ and compare with $\sqrt{\HT}$.
    \item Figure placeholder: classical and quantum scaling curves.
  \end{itemize}
\end{frame}

\begin{frame}{Multiple marked vertices}
  \begin{itemize}
    \item Compare single, independent, and clustered marked sets.
    \item Look for the $\HTp$ penalty in success probability versus time.
    \item Figure placeholder: success probability curves.
  \end{itemize}
\end{frame}

\begin{frame}{Semantic search and decoherence}
  \begin{itemize}
    \item Toy semantic graph: target activation for quantum and classical walks.
    \item Partial measurement sweep: success probability and search time versus $p$.
    \item Is there a nonzero optimal measurement rate?
  \end{itemize}
\end{frame}

\section{Discussion}

\begin{frame}{What the model explains and what it does not}
  \begin{itemize}
    \item Explains: process-level differences in search time and choice probability.
    \item Does not explain: biological implementation or real behavioral data.
    \item Algorithmic advantage is not a biological claim.
    \item Limitations: synthetic graphs, small sizes, simplified noise.
  \end{itemize}
\end{frame}

\begin{frame}{Future work}
  \begin{itemize}
    \item Fit behavioral data with the decision-graph model.
    \item ConceptNet subgraphs and larger semantic networks.
    \item Open-system models beyond projective partial measurement.
    \item Circuit implementation on small graphs.
  \end{itemize}
\end{frame}

\begin{frame}{Conclusion}
  \begin{itemize}
    \item AGJK provides a rigorous search core for a cognitive process model.
    \item Marked-set structure and decoherence are the key knobs.
    \item The project connects quantum algorithms to testable model-level predictions.
  \end{itemize}
\end{frame}

\begin{frame}{References}
  \small
  \begin{itemize}
    \item Ambainis, Gily\'en, Jeffery, Kokainis, arXiv:1903.07493 (2019).
    \item Chen, Ferro, Sornette, PLoS ONE 17(8), e0273551 (2022).
    \item Huang, Busemeyer, ICCM 2022.
    \item Pavone, Faro, Mathematics 13(19), 3171 (2025).
  \end{itemize}
\end{frame}

\end{document}
